\Article{61}
However convenient this method is when the necessary tables are available, we must not forget that it is indirect. So it is worth investigating how much direct methods can achieve. Here we will give what can be drawn from the preceding results, leaving other matters that require deeper considerations for Section VIII. We begin with the simplest case, where $A=1$, that is, where we seek the roots of the congruence $x^n\equiv1\,(\operatorname{mod.}p)$. Taking any primitive root as the base, we must then have $n\operatorname{Ind.}x\equiv0\,(\operatorname{mod.}p-1)$. When $n$ is coprime to $p-1$, this congruence has just one root, namely $\operatorname{Ind.}x\equiv0\,(\operatorname{mod.}p-1)$. Thus \emph{in this case} ${}^{n}\!\surd1\,(\operatorname{mod.}p)$ has a single value, namely $\equiv1$. If, however, the numbers $n,p-1$ have greatest common divisor $\delta$, the complete solution of the congruence $n\operatorname{Ind.}x\equiv0\,(\operatorname{mod.}p-1)$ is $\operatorname{Ind.}x\equiv0\,(\operatorname{mod.}\frac{p-1}{\delta})$ (see Article 29). \emph{In other words,} $\operatorname{Ind.}x$ must be congruent modulo $p-1$ to one of the numbers
\[0,\frac{p-1}{\delta}.\ \frac{2(p-1)}{\delta},\frac{3(p-1)}{\delta},\GaussFourDots\frac{(\delta-1)(p-1)}{\delta}\]
or have $\delta$ values incongruent modulo $p-1$. So $x$ also has $\delta$ different values in this case (incongruent modulo
\SourceLeaf{55}$p$). It follows that the expression ${}^{\delta}\!\surd {1}$ also has $\delta$ different values, whose indices agree exactly with those just given. Thus the expression ${}^{\delta}\!\surd {1}\,(\operatorname{mod.}p)$ is entirely equivalent to ${}^{n}\!\surd {1}\,(\operatorname{mod.}p)$: the congruence $x^\delta\equiv1\,(\operatorname{mod.}p)$ has the same roots as $x^n\equiv1\,(\operatorname{mod.}p)$. The former has lower degree if $\delta$ and $n$ are unequal.

\emph{Example.} ${}^{15}\!\surd {1}\,(\operatorname{mod.}19)$ has three values, since $3$ is the greatest common divisor of $15,18$, and these are also the values of ${}^{3}\!\surd {1}\,(\operatorname{mod.}19)$. They are $1,7,11$.

\Article{62}
This reduction therefore saves us from having to solve any congruences of the form $x^n\equiv1$ except those in which $n$ divides $p-1$. Below we will show that congruences of this form can always be reduced further, although the preceding results are not enough for this. One case can already be settled here, namely $n=2$. The values of ${}^{2}\!\surd {1}$ are clearly $+1$ and $-1$, since there cannot be more than two, and $+1$ and $-1$ are always incongruent unless the modulus is $=2$. In that case it is clear in itself that ${}^{2}\!\surd {1}$ can have only one value. It follows that $+1$ and $-1$ are also the values of ${}^{2m}\!\surd {1}$ when $m$ is coprime to $\frac{p-1}{2}$. This always happens when the modulus is such that $\frac{p-1}{2}$ is a prime number (unless $p-1=2m$, in which case all the numbers $1,2,3\GaussFourDots p-1$ are roots), \emph{for example} when $p=3,5,7,11,23,47,59,83,107$, etc. As a corollary, note that the index of $-1$ is always $\equiv\frac{p-1}{2}\,(\operatorname{mod.}p-1)$, whichever primitive root is taken as the base. For $2\operatorname{Ind.}(-1)\equiv0\,(\operatorname{mod.}p-1)$. Thus $\operatorname{Ind.}(-1)$ is either $\equiv0$ or $\equiv\frac{p-1}{2}\,(\operatorname{mod.}p-1)$. But $0$ is always the index of $+1$, and $+1$ and $-1$ must always have different indices (apart from the case $p=2$, which need not concern us here).

\Article{63}
In Article 60 we showed that the expression ${}^{n}\!\surd {A}\,(\operatorname{mod.}p)$ has $\delta$ different values or none at all, if $\delta$ is the greatest common divisor of $n,p-1$. Just as we have shown that ${}^{n}\!\surd {A}$ and ${}^{\delta}\!\surd {A}$ are equivalent when $A\equiv1$, we will now prove more generally that the expression ${}^{n}\!\surd {A}$ can always be reduced to another, ${}^{\delta}\!\surd {B}$, to which it is equivalent. Denote any value of the former by $x$; then $x^n\equiv A$. Now
\SourceLeaf{56}let $t$ be any value of $\frac{\delta}{n}\,(\operatorname{mod.}p-1)$, which clearly has real values by Article 31. Then $x^{tn}\equiv A^t$, while $x^{tn}\equiv x^\delta$ because $tn\equiv\delta\,(\operatorname{mod.}p-1)$. So $x^\delta\equiv A^t$, and every value of ${}^{n}\!\surd {A}$ is also a value of ${}^{\delta}\!\surd {A^t}$. Whenever ${}^{n}\!\surd {A}$ has real values, it is therefore entirely equivalent to ${}^{\delta}\!\surd {A^t}$, since it has neither other values nor fewer. Yet when ${}^{n}\!\surd {A}$ has no real value, ${}^{\delta}\!\surd {A^t}$ may still have real values.

\emph{Example.} If we seek the values of ${}^{21}\!\surd {2}\,(\operatorname{mod.}31)$, the greatest common divisor of $21$ and $30$ is $3$, and one value of $\frac{3}{21}\,(\operatorname{mod.}30)$ is $3$. So if ${}^{21}\!\surd {2}$ has real values, it is equivalent to ${}^{3}\!\surd {2^3}$, or ${}^{3}\!\surd {8}$. Indeed, the values of the latter expression, namely $2,10.19$, are found to satisfy the former as well.

\Article{64}
To avoid the risk of carrying out this operation for nothing, we need a rule that lets us decide at once whether ${}^{n}\!\surd {A}$ has real values. If a table of indices is available, this is straightforward: Article 60 shows that real values exist if the index of $A$, with any primitive root taken as the base, is divisible by $\delta$, and otherwise do not exist. But this can also be determined without such a table. For if we set the index of $A$ equal to $k$, then, if it is divisible by $\delta$, $\frac{k(p-1)}{\delta}$ is divisible by $p-1$, and conversely. Now the index of $A^{\frac{p-1}{\delta}}$ is $\frac{k(p-1)}{\delta}$. So if ${}^{n}\!\surd {A}\,(\operatorname{mod.}p)$ has real values, $A^{\frac{p-1}{\delta}}$ is congruent to unity; otherwise it is not. Thus in the example of the preceding article we have $2^{10}=1024\equiv1\,(\operatorname{mod.}31)$, from which we conclude that ${}^{21}\!\surd {2}\,(\operatorname{mod.}31)$ has real values.
Similarly, this tells us that ${}^{2}\!\surd {-1}\,(\operatorname{mod.}p)$ always has two real values when $p$ is of the form $4m+1$, but none when $p$ is of the form $4m+3$, because $(-1)^{2m}=1$ and $(-1)^{2m+1}=-1$. This elegant theorem, usually stated as follows: \emph{If $p$ is a prime of the form $4m+1$, a square $aa$ can be found such that $aa+1$ is divisible by $p$; but if $p$ is of the form $4m-1$, no such square exists}, was proved in this way by Euler in \emph{Comm. nov. Acad. Petrop. T.} XVIII, \emph{p.} 112, for the year 1773. He had given another proof much earlier, in \emph{Comm. nov. T.} V, \emph{p.} 5, which appeared in 1760. In an earlier paper, \emph{Comm. nov. T.} IV, \emph{p.} 25, he had not yet
\SourceLeaf{57}completed the matter. Later La Grange also gave a proof of the theorem, \emph{Nouveaux Mém. de l'Ac. de Berlin A.} 1775, \emph{p.} 342. We will give yet another proof in the next section, where this subject will be treated in its own right.

\Article{65}
Having shown how to reduce all expressions ${}^{n}\!\surd {A}\,(\operatorname{mod.}p)$ to ones in which $n$ divides $p-1$, and having obtained a test for whether real values exist, we will examine more closely expressions ${}^{n}\!\surd {A}\,(\operatorname{mod.}p)$ in which $n$ divides $p-1$. First we will show how the individual values of an expression are related; then we will give some methods that often allow one value to be found.

\emph{First,} when $A\equiv1$ and $r$ is one of the $n$ values of ${}^{n}\!\surd {1}\,(\operatorname{mod.}p)$, so that $r^n\equiv1\,(\operatorname{mod.}p)$, all powers of $r$ are also values of that expression. The number of different values among them equals the exponent to which $r$ belongs (Article 48). So if $r$ is a value belonging to the exponent $n$, the powers of $r$, namely $r,r^2,r^3\GaussFourDots r^n$ (the last of which may be replaced by \emph{unity}), include all the values of ${}^{n}\!\surd {1}\,(\operatorname{mod.}p)$. In Section VIII we will discuss more fully the methods available for finding such values belonging to the exponent $n$.

\emph{Second.} When $A$ is incongruent to unity and one value of ${}^{n}\!\surd {A}\,(\operatorname{mod.}p)$, say $z$, is known, the others are obtained as follows. Let the values of ${}^{n}\!\surd {1}$ be

\[1,\ r,\ r^2\ldots r^{n-1}\]

as just shown. Then all the values of ${}^{n}\!\surd {A}$ are

\[z,\ zr,\ zr^2\GaussFourDots zr^{n-1}\]

For all these satisfy $x^n\equiv A$: if any one of them is put $\equiv zr^k$, its power with exponent $n$, namely $z^n r^{nk}$, because of $r^n\equiv1$ and $z^n\equiv A$, is congruent to $A$. Article 23 readily shows that they are all different; and there cannot be more values than these, whose number is $n$, for the expression ${}^{n}\!\surd {A}$. Thus, \emph{for example,} if one value of ${}^{2}\!\surd {A}$ is $z$, the other is $-z$. Finally,
\SourceLeaf{58}we must conclude that all values of ${}^{n}\!\surd {A}$ cannot be found unless all values of ${}^{n}\!\surd {1}$ are also known.

\Article{66}
Our second aim was to show when one value of ${}^{n}\!\surd {A}\,(\operatorname{mod.}p)$ can be found directly, assuming that $n$ divides $p-1$. This happens when some value is congruent to a power of $A$. Since this case occurs quite often, it is worth spending some time on it. Let such a value, \emph{if one exists,} be $z$, so that $z\equiv A^k$ and $A\equiv z^n\,(\operatorname{mod.}p)$. It follows that $A\equiv A^{kn}$. So if we have a number $k$ such that $A\equiv A^{kn}$, then $A^k$ is the value sought. But this condition is equivalent to $1\equiv kn\,(\operatorname{mod.}t)$, where $t$ denotes the exponent to which $A$ belongs (Articles 46, 48). For this congruence to be possible, $n$ must be coprime to $t$. In that case $k\equiv\frac{1}{n}\,(\operatorname{mod.}t)$. But if $t$ and $n$ have a common divisor other than unity, no value $z$ can be congruent to a power of $A$.

\Article{67}
Since this solution requires us to know $t$, let us see how to proceed when that number is unknown. First, it is easy to see that $t$ must divide $\frac{p-1}{n}$ if ${}^{n}\!\surd {A}\,(\operatorname{mod.}p)$ has real values, as we always assume here. For let $y$ be any value. Then both $y^{p-1}\equiv1$ and $y^n\equiv A\,(\operatorname{mod.}p)$ hold. Raising both sides of the latter congruence to the power with exponent $\frac{p-1}{n}$ gives $A^{\frac{p-1}{n}}\equiv1$, so $\frac{p-1}{n}$ is divisible by $t$ (Article 48). Now if $\frac{p-1}{n}$ is coprime to $n$, the congruence $kn\equiv1$ of the preceding article can also be solved modulo $\frac{p-1}{n}$. Any value of $k$ satisfying the congruence for this modulus clearly satisfies it also modulo $t$, which divides $\frac{p-1}{n}$ (Article 5). We have then found what was sought. But if $\frac{p-1}{n}$ is not coprime to $n$, remove from $\frac{p-1}{n}$ all its prime factors that also divide $n$, taking them out of $\frac{p-1}{n}$. This gives a number $\frac{p-1}{nq}$ coprime to $n$, where $q$ denotes the product of all the prime factors removed. If the condition reached in the preceding article holds, namely that $t$ is coprime to $n$, then $t$ is also coprime to $q$, and therefore divides $\frac{p-1}{nq}$ as well. So if we solve $kn\equiv1\,(\operatorname{mod.}\frac{p-1}{nq})$ (which is possible because $n$ is coprime to $\frac{p-1}{nq}$), the value of $k$ will also satisfy the congruence modulo
\SourceLeaf{59}$t$, as required. The whole method consists in finding a number that can take the place of the unknown $t$. But we must remember that when $\frac{p-1}{n}$ is not coprime to $n$, we have assumed the condition of the preceding article to hold. If it fails, all the conclusions will be wrong. If, by following the rules given without regard to this, we find a value for $z$ whose power with exponent $n$ is not congruent to $A$, that tells us that the condition fails and this method cannot be used at all.

\Article{68}
Even in this case, however, the work may often prove useful. It is worth investigating how this false value is related to the true ones. Suppose, then, that $k,z$ have been duly determined, but $z^n$ is not $\equiv A\,(\operatorname{mod.}p)$. If we can determine the values of ${}^{n}\!\surd{\frac{A}{z^n}}\,(\operatorname{mod.}p)$, multiplying each of them by $z$ gives the values of ${}^{n}\!\surd{A}$. For if $v$ is a value of ${}^{n}\!\surd{\frac{A}{z^n}}$, then $(vz)^n\equiv A$. But ${}^{n}\!\surd{\frac{A}{z^n}}$ is simpler than ${}^{n}\!\surd{A}$ in that $\frac{A}{z^n}\,(\operatorname{mod.}p)$ usually belongs to a smaller exponent than $A$. More precisely, if the greatest common divisor of $t,q$ is $d$, then $\frac{A}{z^n}\,(\operatorname{mod.}p)$ belongs to the exponent $d$. This is proved as follows. Substituting the value of $z$ gives $\frac{A}{z^n}\equiv\frac{1}{A^{kn-1}}\,(\operatorname{mod.}p)$. Now $kn-1$ is divisible by $\frac{p-1}{nq}$ (preceding article), while $\frac{p-1}{n}$ is divisible by $t$ (ibid.), or $\frac{p-1}{nd}$ by $\frac{t}{d}$. But $\frac{t}{d}$ is coprime to $\frac{q}{d}$ by hypothesis, so $\frac{p-1}{nd}$ is also divisible by $\frac{tq}{dd}$, or $\frac{p-1}{nq}$ by $\frac{t}{d}$. Thus $kn-1{}^{\prime}$ is also divisible by $\frac{t}{d}$, and $(kn-1)d$ by $t$. Hence $A^{(kn-1)d}\equiv1\,(\operatorname{mod.}p)$. It readily follows that $\frac{A}{z^n}$, raised to the power with exponent $d$, is congruent to unity. It is also easy to prove that $\frac{A}{z^n}$ cannot belong to an exponent smaller than $d$, but since this is not needed for our purpose, we will not dwell on it. We can therefore be certain that $\frac{A}{z^n}\,(\operatorname{mod.}p)$ always belongs to a smaller exponent than $A$, with just one exception: when $t$ divides $q$, so that $d=t$.

But what do we gain from $\frac{A}{z^n}$ belonging to a smaller exponent than $A$? More numbers can occur as $A$ than as $\frac{A}{z^n}$. When several such expressions ${}^{n}\!\surd{A}$ have to be evaluated for the same modulus, this lets us obtain several from the same source. Thus, \emph{for example,} we can always determine at least one value of ${}^{2}\!\surd{A}\,(\operatorname{mod.}29)$,
\SourceLeaf{60}provided that the values of ${}^{2}\!\surd{-1}\,(\operatorname{mod.}29)$ (which are $\pm12$) are known. The preceding article readily shows that one value of such expressions can always be determined directly when $t$ is odd, and that $d$ is $=2$ when $t$ is even. But no number except $-1$ belongs to the exponent $2$.

\emph{Examples.} Find ${}^{3}\!\surd{31}\,(\operatorname{mod.}37)$. Here $p-1=36$, $n=3$, $\frac{p-1}{3}=12$, and so $q=3$. We must therefore have $3k\equiv1\,(\operatorname{mod.}4)$, which is obtained by putting $k=3$. Hence $z\equiv31^3\,(\operatorname{mod.}37)\equiv6$, and indeed we find $6^3\equiv31\,(\operatorname{mod.}37)$. If the values of ${}^{3}\!\surd{1}\,(\operatorname{mod.}37)$ are known, the remaining values of ${}^{3}\!\surd{6}$ can also be determined. The former are $1,10,26$; multiplying $6$ by them gives the remaining values $\equiv23$ and $8$.

If, instead, we seek a value of ${}^{2}\!\surd{3}\,(\operatorname{mod.}37)$, then $n=2$, $\frac{p-1}{n}=18$, and therefore $q=2$. We must have $2k\equiv1\,(\operatorname{mod.}9)$, giving $k\equiv5\,(\operatorname{mod.}9)$. Thus $z\equiv3^5\equiv21\,(\operatorname{mod.}37)$. But $21^2$ is not $\equiv3$, but $\equiv34$. Now $\frac{3}{34}\,(\operatorname{mod.}37)\equiv-1$ and ${}^{2}\!\surd{-1}\,(\operatorname{mod.}37)\equiv\pm6$, giving the true values $\pm6\,.\,21\equiv\pm15$.

This is about as much as we can give here on evaluating such expressions. Direct methods will clearly often turn out to be rather lengthy. But this drawback affects almost all direct methods in number theory, and we did not think it a reason to omit showing what they can achieve here. We should also note that explaining individually the special devices that often suggest themselves to someone experienced is outside our purpose.

\Needspace{10\baselineskip}
\begin{center}\emph{Relations between indices in different systems.}\end{center}

\Article{69}
We now return to the roots we called primitive. We showed that, with any primitive root taken as the base, all numbers whose indices are coprime to $p-1$ are also primitive roots, and no others. This at once gives the number of primitive roots; see Article 53. In general, we may choose whichever primitive root we wish as the base. So here too, as in calculations with logarithms, there can be several systems, so to speak\begingroup\renewcommand{\thefootnote}{*)}\footnote{They differ, however, in that the number of systems for logarithms is infinite, while here it equals the number of primitive roots. Congruent bases clearly generate the same system.}\endgroup.
\SourceLeaf{61}Let us see how they are connected. Let $a,b$ be two primitive roots and $m$ another number. When $a$ is taken as the base, let the index of $b$ be $\equiv\beta$, and that of $m$ be $\equiv\mu\,(\operatorname{mod.}p-1)$. When $b$ is taken as the base, let the index of $a$ be $\equiv\alpha$, and that of $m$ be $\equiv\nu\,(\operatorname{mod.}p-1)$. Then $\alpha\beta\equiv1\,(\operatorname{mod.}p-1)$. For $a^\beta\equiv b$, so $a^{\alpha\beta}\equiv b^\alpha\equiv a\,(\operatorname{mod.}p)$ by hypothesis, and hence $\alpha\beta\equiv1\,(\operatorname{mod.}p-1)$. Similar reasoning gives $\nu\equiv\alpha\mu$ and $\mu\equiv\beta\nu\,(\operatorname{mod.}p-1)$. So a table of indices constructed for the base $a$ can easily be converted into one with base $b$. For if, for the base $a$, the index of $b$ is $\equiv\beta$, then for the base $b$ the index of $a$ is $\equiv\frac{1}{\beta}\,(\operatorname{mod.}p-1)$. Multiplying all the indices in the table by this number gives all the indices for the base $b$.

\Article{70}
Although a given number may have several indices when different primitive roots are taken as the base, they all agree in having the same greatest common divisor with $p-1$. For suppose the index of the given number for the base $a$ is $m$, and for the base $b$ is $n$. Suppose their greatest common divisors with $p-1$, namely $\mu,\nu$, are unequal. One is then larger, say $\mu>\nu$, so $\mu$ does not divide $n$. But if the index of $a$ when $b$ is taken as the base is denoted by $\alpha$, then, by the preceding article, $n\equiv\alpha m\,(\operatorname{mod.}p-1)$. Thus $\mu$ also divides $n$, \emph{contrary to the assumption.}

That this greatest common divisor of the indices of a given number and $p-1$ is independent of the base also follows from its being equal to $\frac{p-1}{t}$, where $t$ is the exponent to which the number whose indices we are considering belongs. For if the index for any base is $k$, then $t$ is the smallest number which, multiplied by $k$, gives a multiple of $p-1$, excluding zero (see Articles 48, 58). In other words, it is the smallest value of $\frac{0}{k}\,(\operatorname{mod.}p-1)$ other than zero. That this equals the greatest common divisor of $k$ and $p-1$ follows readily from Article 29.

\Article{71}
It is also easy to prove that the base can always be chosen so that a number belonging to the exponent $t$ has any given index whose greatest common divisor with $p-1$ is $=\frac{p-1}{t}$. Denote this divisor briefly by $d$, and let the desired index be $\equiv dm$. When any
\SourceLeaf{62}primitive root $a$ is taken as the base, let the index of the given number be $\equiv dn$. Then $m,n$ are coprime to $\frac{p-1}{d}$, or $t$. Now if $\varepsilon$ is a value of $\frac{dn}{dm}\,(\operatorname{mod.}p-1)$ and is also coprime to $p-1$, then $a^\varepsilon$ is a primitive root. Taking it as the base gives the number in question the index $dm$, as required (for $a^{\varepsilon dm}\equiv a^{dn}\equiv$ the given number). That $\frac{dn}{dm}\,(\operatorname{mod.}p-1)$ has values coprime to $p-1$ is proved as follows. This expression is equivalent to $\frac{n}{m}\,(\operatorname{mod.}\frac{p-1}{d})$, or $\frac{n}{m}\,(\operatorname{mod.}t)$; see Article 31, 2. All its values are coprime to $t$. For if a value $e$ had a common divisor with $t$, this divisor would also have to divide $me$, and therefore $n$, to which $me$ is congruent modulo $t$. This contradicts the hypothesis that $n$ is coprime to $t$. So when every prime divisor of $p-1$ also divides $t$, \emph{all} values of $\frac{n}{m}\,(\operatorname{mod.}t)$ are coprime to $p-1$, and their number is $=d$. But if $p-1$ has other prime divisors $f,g,h$, etc., that do not divide $t$, choose any value, writing $\frac{n}{m}\,(\operatorname{mod.}t)\equiv e$. Since $t,f,g,h$, etc., are all pairwise coprime, a number $\varepsilon$ can be found that is congruent modulo $t$ to $e$ and, modulo $f,g,h$, etc., to any numbers respectively coprime to them (Article 32). Such a number is divisible by no prime factor of $p-1$, and is therefore coprime to $p-1$, as required. Finally, the theory of combinations readily gives $=\frac{p-1\,.\,f-1\,.\,g-1\,.\,h-1.\ \mathrm{etc.}}{t\,.\,f\,.\,g\,.\,h\,.\,\mathrm{etc.}}$ for the number of such values. But to keep this digression from growing too long, we omit the proof, since it is not particularly needed for our purpose.

\begin{center}\emph{Bases suited to particular uses.}\end{center}

\Article{72}
Although the choice of a primitive root as the base is generally entirely arbitrary, some bases can sometimes offer particular advantages over others. In Table I we always took $10$ as the base when it was a primitive root. Otherwise we always chose the base so that the index of $10$ was as small as possible, \emph{that is,} $=\frac{p-1}{t}$, where $t$ denotes the exponent to which $10$ belongs. In Section VI we will show what this gains us, when the same table is put to other uses as well. But since some choice may still remain here, as the preceding article shows, we adopted a fixed rule: among all primitive roots meeting the requirement, we always chose the \emph{smallest} as the base. Thus for $p=73$, where $t=8$ and $d=9$, $a^\varepsilon$ has $\frac{72\,.\,2}{8\,.\,3}$, \emph{that is,} $6$ values, namely $5,14,20,28,39,40$. So we took the smallest, $5$, as the base.
